\section*{EPILOGUE}

Me voilà au terme de cette brève et très incomplète histoire de Meyrueis reconstituée à partir
de quelques documents que j’ai pu retrouver et rassembler, de notices et conversations
familiales ou amicales, d’articles ou de notes empruntées à des journaux ou revues diverses.

On y discerne, mais ce n’est guère original, une évolution des modes de pensée et des modes
de vie accélérée par plusieurs ébranlements.

\textbf{Le premier fut l’irruption de la Réforme dès 1550,}

bousculant les structures de la société médiévale, remettant en question la toute puissance
religieuse, politique et économique de l’Eglise Catholique Romaine, permettant au peuple de
comprendre dans sa langue, le message subversif de la Bible.

\textbf{Le deuxième le courant philosophique des Lumières 1650-1700,}

inspiré par les œuvres de Fontenelle et surtout de Bayle, protestant réfugié à Genève puis
Rotterdam au temps des persécutions, relayé par Montesquieu, Diderot, Voltaire et tous les
Encyclopédistes qui font toute confiance à la raison humaine ! \emph{« La tolérance est une vertu
primordiale et la raison devrait conduire l’humanité toute entière à la sagesse et à la perfection ».}
(Le petit Robert p. 1313). On pouvait trouver certains de ces livres dans la bibliothèque du château de Roquedols.
Toute l’effervescence de ce courant débouche sur la Révolution, et, après les années terribles
1789 -- 1794, une amélioration, de 1800 à 1914, de la vie des petits paysans et artisans de nos
campagnes, (les 85\% de la population française).

\textbf{Le troisième ce fut la Grande Guerre 1914-1918.}

Des millions d’hommes jeunes et en pleine force mobilisés : les femmes prirent la relève dans
bien des domaines « réservés » aux hommes, et s’en tirèrent plutôt bien. Beaucoup ne
revinrent pas et il n’est que de s’arrêter devant le monument aux morts de Meyrueis, avec ses
84 noms auxquels on peut ajouter les 6 de Gatuzìère, pour réaliser les conséquences et le
désastre : des femmes abandonnées entre les deux guerres, de petites entreprises également :
chapelleries, filatures, tissages...

\textbf{Le quatrième la Seconde Guerre Mondiale 1939-1945.}

Elle toucha Meyrueis de manière différente que celle de 14-18. Moins de morts certes mais
l’arrivée à Meyrueis entre 1939 et 1944 de nombreuses familles réfugiées, surtout en 1939-
1941, craignant les combats ou chassées par les lois visant les juifs, les étrangers et les
apatrides qui se cachèrent dans nos hameaux et villages.

De plus l’installation d’un groupement des Chantiers de la Jeunesse, (1600 jeunes à Meyrueis
et dans les environs !), organisation paramilitaire de Vichy, bouleversa la vie bien réglée
d’avant guerre au point qu’en 1941 il y eut une disette de pommes de terre à Meyrueis !!
Les restrictions alimentaires eurent pour conséquence, avec la distribution de cartes
d’alimentation, de vêtements ou de chaussures, l’émergence de ce que l’on a appelé le marché
noir : (certains vendaient à des prix parfois très élevés des marchandises rationnées). Peut-on
penser que ce fut seulement un phénomène marginal ? Peut-être. Ce qui est sûr, c’est qu’il y
eut, à Meyrueis, des familles intègres, et accueillantes pour les pourchassés de toutes sortes où
le pain et le lard étaient partagés...Mais après 1945 la vie ne fut plus tout à fait la même....
